% Part: turing-machines
% Chapter: machines-computations
% Section: combining-machines

\documentclass[../../../include/open-logic-section]{subfiles}

\begin{document}

\olfileid{tur}{mac}{cmb}
\olsection{ట్యూరింగ్ యంత్రాలను కలపడం}

\begin{explain}
ఇంతవరకు మనం చూసిన ట్యూరింగ్ యంత్రాల ఉదాహరణలు చాలా సరళమైనవి. కానీ నిజానికి,
ఏ ఆధునిక ప్రోగ్రామింగ్ భాషతో పరిష్కరించగల ఏ సమస్యనైనా ట్యూరింగ్ యంత్రాలతో
కూడా పరిష్కరించవచ్చు. మరింత సంక్లిష్టమైన ట్యూరింగ్ యంత్రాలను నిర్మించడానికి,
యంత్రాలను కలపగలమని మనకు మనమే నమ్మకం కలిగించుకోవడం ముఖ్యం. అప్పుడు ప్రక్రియను
సరళమైన భాగాలుగా విడగొట్టి, మరింత సంక్లిష్ట సమస్యలను పరిష్కరించే యంత్రాలను
నిర్మించవచ్చు. సంక్లిష్ట సమస్యను దాని భాగాలుగా విభజించడానికి సహజమైన పద్ధతిని
కనుగొనగలిగితే, సమస్యను అనేక దశల్లో ఎదుర్కొనవచ్చు: అనేక సరళ ట్యూరింగ్
యంత్రాలను నిర్మించి, వాటిని ఒకే సమస్య-పరిష్కార యంత్రంగా కలపవచ్చు. తరువాతి
విభాగంలో ఆగే సమస్యను ఎదుర్కొనేటప్పుడు ఈ విషయం ప్రత్యేకంగా ముఖ్యం.

ట్యూరింగ్ యంత్రాలు $M = \tuple{Q, \Sigma, q_0, \delta}$,
~$M' = \tuple{Q', \Sigma', q_0', \delta'}$లను ఎలా కలుపుతాం? $M$ ఆగిన
తరువాతి టేపు స్థితివిన్యాసాన్ని ఇప్పుడు యంత్రం~$M'$ నడకకు ఇన్‌పుట్
స్థితివిన్యాసంగా ఉపయోగిస్తాం. ఇలా చేసే ఒకే ట్యూరింగ్ యంత్రం $M \frown M'$ను
పొందడానికి కింది విధంగా చేయండి:
\begin{enumerate}
    \item $M$, $M'$లకు ఏ స్థితీ ఉమ్మడిగా లేకుండా ($Q \cap Q' = \emptyset$)
    $M'$ యొక్క అన్ని స్థితులు~$Q'$కు మళ్లీ సంఖ్యలు (లేదా పేర్లు) ఇవ్వండి.
    \item $M \frown M'$ యొక్క స్థితులు $Q \cup Q'$.
    \item టేపు వర్ణమాల $\Sigma \cup \Sigma'$.
    \item ప్రారంభ స్థితి~$q_0$.
    \item సంక్రమణ ప్రమేయం కింది విధంగా ఇచ్చిన ప్రమేయం $\delta''$:
    \[\delta''(q,\sigma) =
    \begin{cases}
      \delta(q,\sigma) & \text{$q \in Q$, $\delta(q,\sigma)$ నిర్వచితమైతే}\\
      \delta'(q,\sigma) & \text{$q \in Q'$ అయితే}\\
      \tuple{q_0', \sigma, \TMstay} & \text{$q \in Q$, అలాగే
      $\delta(q,\sigma)$ నిర్వచితం కాకపోతే}
    \end{cases}\]
\end{enumerate}
\sourcecorrection{OLTETURMAC-007}{ఆంగ్ల మూలంలోని $\delta''$ నిర్వచనంలో మొదటి సందర్భం $q\in Q$ అని మాత్రమే ఉండి, మూడవ ``$q\in Q$, $\delta(q,\sigma)$ నిర్వచితం కాదు'' సందర్భంతో అతివ్యాప్తి చెందింది. తరువాతి వివరణ నిర్దేశించినట్లుగా, మొదటి సందర్భానికి $\delta(q,\sigma)$ నిర్వచితమనే షరతును చేర్చి సందర్భాలను విడదీశాం.}
ఫలితంగా వచ్చే యంత్రం స్థితి $q \in Q$లో ఉన్నప్పుడు $M$ యొక్క నిర్దేశాలను,
స్థితి~$q \in Q'$లో ఉన్నప్పుడు $M'$ యొక్క నిర్దేశాలను ఉపయోగిస్తుంది. అది
స్థితి $q \in Q$లో ఉండి, $M$కు సంక్రమణ లేని సంకేతం~$\sigma$ను చదువుతుంటే
(అంటే $M$ ఆగిపోయి ఉండేది), $M'$ యొక్క ప్రారంభ స్థితిలోకి వెళుతుంది (టేపు
విషయాన్ని, శీర్షం స్థానాన్ని యథాతథంగా ఉంచుతుంది).

యంత్రం~$M$ క్రమశిక్షితమైనది కాకపోతే, $M$ ఆగినప్పుడు టేపు శీర్షం ఎక్కడ ఉందో
మనకు తెలియదని గమనించండి; అందువల్ల $M$ యొక్క ఆగే స్థితివిన్యాసంలో శీర్షం
గడి~$1$ను చదువుతూ ఉండనవసరం లేదు. యంత్రాలను కలిపేటప్పుడు దీన్ని గుర్తుంచుకోవడం
ముఖ్యం.
\end{explain}

\begin{ex}
\emph{యంత్రాలను కలపడం:} పొడవులు $n$, $m$ అయిన రెండు $\TMstroke$ల దిమ్మలతో
కూడిన ఇన్‌పుట్‌పై ప్రారంభించినప్పుడు, టేపుపై $2(m+n)$ $\TMstroke$ల ఒకే
దిమ్మతో ఆగే యంత్రాన్ని రూపొందిస్తాం. ఈ యంత్రాన్ని నిర్మించడానికి, మనకు ఇప్పటికే
తెలిసిన కూడిక యంత్రం, ద్విగుణక యంత్రం అనే రెండు యంత్రాలను కలపవచ్చు. కూడిక
యంత్రానికి స్థితి రేఖాచిత్రం గీయడంతో మొదలుపెడదాం.
\[
\begin{tikzpicture}[->,>=stealth',shorten >=1pt,auto,node distance=2.8cm,
                    semithick]
  \tikzstyle{every state}=[fill=none,draw=black,text=black]

  \node[initial,state] (A)              {$q_0$};
  \node[state]         (B) [right of=A] {$q_1$};
  \node[state]         (C) [right of=B] {$q_2$};

  \path (A) edge node {\TMtrans{\TMblank}{\TMstroke}{\TMstay}} (B)
            edge [loop above] node {\TMtrans{\TMstroke}{\TMstroke}{\TMright}} (A)
        (B) edge [loop above] node {\TMtrans{\TMstroke}{\TMstroke}{\TMright}} (B)
            edge node {\TMtrans{\TMblank}{\TMblank}{\TMleft}} (C)
        (C) edge [loop above] node {\TMtrans{\TMstroke}{\TMblank}{\TMstay}} (C);
\end{tikzpicture}
\]
స్థితి~$q_2$లో ఆగడానికి బదులుగా, అవుట్‌పుట్‌ను ద్విగుణీకరించడానికి పనిని
కొనసాగించాలి. ద్విగుణక యంత్రం ఇన్‌పుట్‌లోని మొదటి గీతను చెరిపి, వేరుగా ఉన్న
అవుట్‌పుట్‌లో రెండు గీతలను వ్రాస్తుందని గుర్తుచేసుకోండి. టేపు శీర్షం కూడిక
యంత్రపు అవుట్‌పుట్‌లోని మొదటి గీతను చదువుతోందని నిర్ధారించడానికి ఒక నిర్దేశాన్ని
చేర్చుదాం.
\[
\begin{tikzpicture}[->,>=stealth',shorten >=1pt,auto,node distance=2.8cm,
                    semithick]
  \tikzstyle{every state}=[fill=none,draw=black,text=black]

  \node[initial,state] (A)              {$q_0$};
  \node[state]         (B) [right of=A] {$q_1$};
  \node[state]         (C) [right of=B] {$q_2$};
  \node[state]         (D) [below left of=C] {$q_3$};
  \node[state]         (E) [below left of=D] {$q_4$};

  \path (A) edge node {\TMtrans{\TMblank}{\TMstroke}{\TMstay}} (B)
            edge [loop above] node {\TMtrans{\TMstroke}{\TMstroke}{\TMright}} (A)
        (B) edge [loop above] node {\TMtrans{\TMstroke}{\TMstroke}{\TMright}} (B)
            edge node {\TMtrans{\TMblank}{\TMblank}{\TMleft}} (C)
        (C) edge node {\TMtrans{\TMstroke}{\TMblank}{\TMleft}} (D)
        (D) edge [loop left] node {\TMtrans{\TMstroke}{\TMstroke}{\TMleft}} (D)
            edge node[left, xshift=-2mm] {\TMtrans{\TMendtape}{\TMendtape}{\TMright}} (E);
\end{tikzpicture}
\]
ఇప్పుడు ఇన్‌పుట్‌ను ద్విగుణీకరించడం సులభం---స్థితి~$q_4$కు ద్విగుణక యంత్రాన్ని
కలిపితే చాలు. దీనికోసం ద్విగుణక యంత్రపు స్థితులకు మళ్లీ పేర్లు ఇచ్చి, అవి
~$q_0$ వద్ద కాక~$q_4$ వద్ద మొదలయ్యేలా చేయాలి---ఈ విధంగా రెండు ప్రారంభ
స్థితులు ఏర్పడవు. చివరి రేఖాచిత్రం \olref{fig:combined}లో ఉన్నట్లుగా ఉండాలి.
\begin{figure}
\[
\begin{tikzpicture}[->,>=stealth',shorten >=1pt,auto,node distance=2.8cm,
                    semithick]
  \tikzstyle{every state}=[fill=none,draw=black,text=black]
  \node[initial,state] (A)              {$q_0$};
  \node[state]         (B) [right of=A] {$q_1$};
  \node[state]         (C) [right of=B] {$q_2$};
  \node[state]         (D) [below left of=C] {$q_3$};
  \node[state]         (E) [below left of=D] {$q_4$};
  \node[state]         (2) [right of=E] {$q_5$};
  \node[state]         (3) [right of=2] {$q_6$};
  \node[state]         (4) [below of=3] {$q_7$};
  \node[state]         (5) [left of=4]  {$q_8$};
  \node[state]         (6) [left of=5]  {$q_9$};

  \path (A) edge node {\TMtrans{\TMblank}{\TMstroke}{\TMstay}} (B)
            edge [loop above] node {\TMtrans{\TMstroke}{\TMstroke}{\TMright}} (A)
        (B) edge [loop above] node {\TMtrans{\TMstroke}{\TMstroke}{\TMright}} (B)
            edge node {\TMtrans{\TMblank}{\TMblank}{\TMleft}} (C)
        (C) edge node {\TMtrans{\TMstroke}{\TMblank}{\TMleft}} (D)
        (D) edge [loop left] node {\TMtrans{\TMstroke}{\TMstroke}{\TMleft}} (D)
            edge node[left, xshift=-2mm] {\TMtrans{\TMendtape}{\TMendtape}{\TMright}} (E)
    (E) edge node {\TMtrans{\TMstroke}{\TMblank}{\TMright}} (2)
    (2) edge [loop above] node {\TMtrans{\TMstroke}{\TMstroke}{\TMright}} (2)
      edge node {\TMtrans{\TMblank}{\TMblank}{\TMright}} (3)
    (3) edge [loop above] node {\TMtrans{\TMstroke}{\TMstroke}{\TMright}} (3)
        edge node {\TMtrans{\TMblank}{\TMstroke}{\TMright}} (4)
    (4) edge [loop below] node {\TMtrans{\TMblank}{\TMstroke}{\TMleft}} (4)
        edge node {\TMtrans{\TMstroke}{\TMstroke}{\TMleft}} (5)
    (5) edge [loop below]  node {\TMtrans{\TMstroke}{\TMstroke}{\TMleft}} (5)
        edge              node {\TMtrans{\TMblank}{\TMblank}{\TMleft}} (6)
    (6) edge [loop below] node {\TMtrans{\TMstroke}{\TMstroke}{\TMleft}} (6)
        edge              node {\TMtrans{\TMblank}{\TMblank}{\TMright}} (E);
\end{tikzpicture}
\]
\caption{కూడిక, ద్విగుణక యంత్రాలను కలపడం}
\ollabel{fig:combined}
\end{figure}
\sourcecorrection{OLTETURMAC-008}{ఆంగ్ల మూలం కూడిక యంత్రాన్ని ఈ ఉదాహరణలో మూడుసార్లు పునర్వాడుతూ, ప్రతి సారి స్థితి $q_0$ వద్ద $\TMstroke$ బాణాన్ని స్వీయ-చక్రంగా అలంకరించినా లక్ష్యాన్ని $q_1$గా ఇచ్చింది. కూడిక దశ మొదటి ఇన్‌పుట్ దిమ్మను పూర్తిగా దాటేందుకు మూడు లక్ష్యాలనూ $q_0$గా సరిచేశాం.}
\end{ex}

\begin{prop}
    $M$, $M'$లు క్రమశిక్షితమైనవై, వరుసగా ప్రమేయాలు $f\colon
    \Nat^k \to \Nat$, $f'\colon \Nat \to \Nat$లను గణిస్తే,
    $M \frown M'$ క్రమశిక్షితమైనది, అలాగే~$\comp{f}{f'}$ను గణిస్తుంది.
\end{prop}

\begin{proof}
    $M$ క్రమశిక్షితమైనది కనుక, అవుట్‌పుట్~$f(n_1,\dots,n_k) = m$తో
    ఆగినప్పుడు శీర్షం గడి~$1$ను చదువుతూ ఉంటుంది. ఇప్పుడు~$M'$ యొక్క ప్రారంభ
    స్థితిలోకి వెళితే, యంత్రం అవుట్‌పుట్ $f'(m)$తో, మళ్లీ గడి~$1$ను చదువుతూ
    ఆగుతుంది. \olref[dis]{defn:disciplined}లోని ఇతర షరతులు కూడా నెరవేరుతాయి.
\end{proof}

\begin{prob}
    \cref{tur:mac:dis:prob:disc-succ}లోని యంత్రం~$M$ను తీసుకుని
    $M \frown M$ను నిర్మించడం ద్వారా $f(x) = x+2$ను గణించే క్రమశిక్షిత
    ట్యూరింగ్ యంత్రాన్ని ఇవ్వండి.
\end{prob}
\end{document}
